\documentclass[10pt,a4paper]{article}
\usepackage[utf8]{inputenc}
\usepackage[italian]{babel}
\usepackage{amsmath}
\usepackage{amsfonts}
\usepackage{amssymb}
\usepackage{geometry}
\usepackage{multicol}
\usepackage{xcolor}
\usepackage{tcolorbox}

 \usepackage{hyperref}
\hypersetup{colorlinks=true, linkcolor=blue, urlcolor=blue}

\geometry{margin=1.5cm}

\title{\textbf{Algebra Lineare: Guida risoluzione esercizi}}
\author{Alessio Avarappattu}

\begin{document}

\maketitle
\section*{}

Questa guida segue la linea didattica del corso di \textbf{Algebra Lineare} del Prof. Giovanni Gaiffi (Ingegneria Informatica, @Unipi).

\medskip

\textbf{Prerequisiti}
\\ I concetti teorici fondamentali sono dati per assodati. Per chi volesse rivedere i concetti principali (schema riassuntivo), gran parte della trattazione teorica è stata pubblicata ed è consultabile liberamente su \href{https://uni.alessioavara.it/algebra-lineare/teoria}{\textbf{uni.alessioavara.it}}.
Questo documento non sostituisce lo studio della teoria, ha il solo obiettivo di metterla in pratica.

\medskip


\vspace{0.5cm}
\hrule
\vspace{1cm}
\section*{1. Calcolo Dimensioni e Basi: Nucleo (Kernel) e Immagine}
\textit{Procedure standard per determinare le proprietà fondamentali di trasformazioni lineari espresse sia tramite matrici esplicite che tramite definizioni funzionali.}

\hrule
\vspace{0.4cm}

\begin{multicols}{2}

\subsection*{1.1 Calcolo con MATRICI}
L'approccio matriciale si basa sulla riduzione a scala (Algoritmo di Gauss).

\medskip
\textbf{Procedura:}
\begin{enumerate}
    \item Scrivere la matrice $A$ associata.
    \item Ridurre $A$ in una forma a scala.
    \item \textbf{Immagine ($Im$):} La dimensione è pari al numero di pivot. Una base è data dalle colonne della matrice \textbf{originale} $A$ in corrispondenza dei pivot.
    \item \textbf{Nucleo ($Ker$):} Risolvere il sistema omogeneo $Ax = 0$. La dimensione è $n - r$ (colonne - pivot).
\end{enumerate}

\textbf{Esempio Svolto:}
$A = \begin{pmatrix} 1 & 2 & 0 \\ 1 & 2 & 1 \\ 0 & 0 & 1 \end{pmatrix} \xrightarrow{\text{Gauss}} \begin{pmatrix} \mathbf{1} & 2 & 0 \\ 0 & 0 & \mathbf{1} \\ 0 & 0 & 0 \end{pmatrix}$

\begin{itemize}
    \item \textbf{Dim Im:} 2 (due pivot).
    \item \textbf{Base Im:} $\{ (1, 1, 0), (0, 1, 1) \}$.
    \item \textbf{Dim Ker:} $3 - 2 = 1$.
    \item \textbf{Base Ker:} Risolvendo il sistema, ponendo $y=1$ otteniamo $\{ (-2, 1, 0) \}$.
\end{itemize}

\columnbreak

\subsection*{1.2 Calcolo con APPLICAZIONI}
Nelle applicazioni (polinomi, funzioni), si costruisce la matrice rappresentativa rispetto alle basi canoniche.

\medskip
\textbf{Procedura:}
\begin{enumerate}
    \item Scegliere le basi canoniche per dominio e codominio.
    \item Calcolare l'immagine di ogni vettore della base del dominio.
    \item Inserire i vettori risultanti in \textbf{colonna} per formare la matrice $M$.
    \item Procedere con Gauss su $M$.
\end{enumerate}

\textbf{Esempio Svolto ($T: P_2 \to P_1, T(p)=p'$):}
Basi: $\mathcal{B}_{in} = \{1, x, x^2\}$, $\mathcal{B}_{out} = \{1, x\}$.

\begin{itemize}
    \item $T(1) = 0 \to \text{colonna } (0, 0)^T$
    \item $T(x) = 1 \to \text{colonna } (1, 0)^T$
    \item $T(x^2) = 2x \to \text{colonna } (0, 2)^T$
\end{itemize}

Matrice $M = \begin{pmatrix} 0 & 1 & 0 \\ 0 & 0 & 2 \end{pmatrix}$
\begin{itemize}
    \item \textbf{Dim Im:} 2. Base: $\{1, x\}$ (suriettiva).
    \item \textbf{Dim Ker:} $3 - 2 = 1$. Base: $\{1\}$ (costanti).
\end{itemize}

\end{multicols}


\vspace{0.3cm}
\hrule
\vspace{0.5cm}

\section*{2. Iniettività e Suriettività: Criteri Rapidi}
Sia $f: V \to W$ un'applicazione lineare e $A$ la matrice associata $m \times n$ (dove $n = \dim(V)$ e $m = \dim(W)$).

\begin{itemize}
    \item \textbf{Iniettività:} $f$ è iniettiva se e solo se trasforma vettori distinti in immagini distinte.
    \begin{itemize}
        \item \textbf{Condizione Nucleo:} $\ker(f) = \{0\}$.
        \item \textbf{Condizione Rango:} $rg(A) = n$ (il rango è massimo, pari al numero di colonne).
        \item \textbf{Nota Pratica:} Se ci sono colonne senza pivot, \textbf{non} è iniettiva.
        \item \textbf{Matrici Alte ($m > n$):} Possono essere iniettive, mai suriettive.
    \end{itemize}
    
    \item \textbf{Suriettività:} $f$ è suriettiva se l'Immagine coincide con il Codominio.
    \begin{itemize}
        \item \textbf{Condizione Immagine:} $\dim(\text{Im } f) = m$.
        \item \textbf{Condizione Rango:} $rg(A) = m$ (il rango è pari al numero di righe).
        \item \textbf{Nota Pratica:} Se dopo Gauss c'è una riga nulla (e $m=n$), \textbf{non} è suriettiva.
        \item \textbf{Matrici Larghe ($n > m$):} Possono essere suriettive, mai iniettive.
    \end{itemize}
\end{itemize}

\textbf{Sintesi per Matrici Quadrate ($n \times n$):}
Se $\det(A) \neq 0$, l'applicazione è un \textbf{Isomorfismo} (Biettiva).



\vspace{0.5cm}
\hrule
\vspace{0.5cm}

\section*{3. Determinazione di Immagini tramite Linearità}
\textit{Metodo per calcolare l'immagine di un vettore generico conoscendo le immagini di una base o di un set di vettori.}

\textbf{Esempio:} Dati $f(e_1+e_2+e_3)=e_2, f(e_1-e_3)=e_1, f(e_2-e_1)=e_2$. Trovare $f(e_1)$.

\textbf{Svolgimento:}
Poniamo $v_1=f(e_1), v_2=f(e_2), v_3=f(e_3)$. Impostiamo il sistema vettoriale:
\begin{enumerate}
    \item $v_1 + v_2 + v_3 = e_2$
    \item $v_1 - v_3 = e_1 \implies v_3 = v_1 - e_1$
    \item $v_2 - v_1 = e_2 \implies v_2 = v_1 + e_2$
\end{enumerate}

Sostituendo le espressioni (2) e (3) nella (1):
\[ v_1 + (v_1 + e_2) + (v_1 - e_1) = e_2 \]
\[ 3v_1 - e_1 = 0 \implies 3v_1 = e_1 \]
\[ \mathbf{f(e_1) = \frac{1}{3}e_1} \]

\vspace{0.5cm}
\hrule
\vspace{0.5cm}

\section*{4. Verifica della Linearità}
Analisi di mappe $T: X \to X$ per determinare la linearità al variare di parametri.

\begin{tcolorbox}[colback=blue!5!white,colframe=blue!75!black,title=Nota Fondamentale sul Calcolo]
Durante la verifica dell'omogeneità $T(kf)$, lo scalare $k$ è una \textbf{costante reale}. 
Secondo le regole di derivazione, la costante esce dall'operatore derivata:
\[ (k \cdot f(x))' = k \cdot f'(x) \]
\textbf{Attenzione:} Non scrivere $k'$ (derivata di $k$ rispetto a $x$), poiché essendo costante varrebbe 0.
\end{tcolorbox}

\subsection*{4.1 Algoritmo di Verifica}
\begin{enumerate}
    \item \textbf{Test dello Zero:} Calcola $T(0)$. Se $T(0) \neq 0$, allora \textbf{NON LINEARE}.
    \item \textbf{Test della Costante ($k$):} Calcola $T(kf)$ sostituendo ogni $f$ con $kf$. Confronta con $k \cdot T(f)$. Se compaiono potenze ($k^2$) o valori assoluti ($|k|$), \textbf{NON LINEARE}.
    \item \textbf{Test della Somma:} Verifica $T(f+g) = T(f) + T(g)$.
\end{enumerate}

\subsection*{4.2 Esercizi Svolti}

\textbf{Esercizio 1: $T(f)(x) = |f'(x)| + a f(x)$}
\begin{itemize}
    \item \textbf{Verifica Omogeneità:} Calcoliamo $T(kf) = |(kf)'| + a(kf)$.
    \item Poiché $k$ è costante: $|k f'| + akf = |k| \cdot |f'| + akf$.
    \item Confrontiamo con $k \cdot T(f) = k |f'| + akf$.
    \item \textbf{Analisi:} Se $k < 0$, allora $|k| \neq k$. L'uguaglianza non regge.
    \item \textbf{Soluzione:} Mai lineare per alcun valore di $a$.
\end{itemize}

\textbf{Esercizio 2: $T(f)(x) = f(x)^a + b$}
\begin{itemize}
    \item \textbf{Test dello Zero:} $T(0) = 0^a + b = b$. Per essere lineare, deve essere $b=0$.
    \item \textbf{Test Omogeneità (con $b=0$):} $T(kf) = (kf)^a = k^a f^a$.
    \item Confrontiamo con $k T(f) = k f^a$. L'uguaglianza $k^a = k$ vale per ogni $k$ solo se $a=1$.
    \item \textbf{Soluzione:} Lineare se e solo se $a=1$ e $b=0$.
\end{itemize}

\textbf{Esercizio 3: $T(f)(x) = f(x^2) + a x f'(x)$}
\begin{itemize}
    \item \textbf{Test Omogeneità:} $T(kf) = (kf)(x^2) + ax(kf)'(x)$.
    \item Applicando la proprietà della costante: $k f(x^2) + ax(k f'(x))$.
    \item Raccogliamo $k$: $k [ f(x^2) + ax f'(x) ]$.
    \item Questo coincide esattamente con $k \cdot T(f)$. Le operazioni sull'argomento ($x^2$) non influenzano la linearità di $f$.
    \item \textbf{Soluzione:} Lineare per ogni $a \in \mathbb{R}$.
\end{itemize}

\textbf{Esercizio 4: $T(f) = \int_{0}^{1} [f(t) + a]^2 dt$}
\begin{itemize}
    \item Sviluppando il quadrato: $T(f) = \int (f^2 + 2af + a^2) dt$.
    \item Il termine $f^2$ comporta che $T(kf)$ conterrà un termine $k^2 f^2$, violando la condizione di omogeneità ($k^2 \neq k$).
    \item \textbf{Soluzione:} Mai lineare.
\end{itemize}

\textbf{Esercizio 5: $T(f)(x) = a f'(x+1) + b f(x) + c$}
\begin{itemize}
    \item \textbf{Test dello Zero:} $T(0) = a(0) + b(0) + c = c$. Condizione necessaria: $c=0$.
    \item \textbf{Test Omogeneità (con $c=0$):} $T(kf) = a(kf)'(x+1) + b(kf)(x)$.
    \item Estraiamo la costante $k$: $k [ a f'(x+1) + b f(x) ]$.
    \item L'espressione coincide con $k T(f)$.
    \item \textbf{Soluzione:} Lineare se $c=0$, per ogni $a, b \in \mathbb{R}$.
\end{itemize}
\section*{5. Sottospazi Vettoriali: Intersezione e Somma}
\textit{Analisi delle dimensioni, rappresentazioni cartesiane vs parametriche e Teorema di Grassmann.}

\hrule
\vspace{0.4cm}

\subsection*{5.1 Rappresentazioni dei Sottospazi}
Un sottospazio $U \subseteq \mathbb{R}^n$ può essere descritto in due modi:
\begin{enumerate}
    \item \textbf{Rappresentazione Parametrica (Span):}
    $U = \text{span}\{v_1, v_2, \dots, v_k\}$.
    È una matrice che genera vettori. La dimensione è il numero di vettori linearmente indipendenti.
    
    \item \textbf{Rappresentazione Cartesiana (Equazioni):}
    $U = \{ x \in \mathbb{R}^n \mid Ax = 0 \}$.
    È un'unica equazione/sistema composto da equazioni omogenee.
\end{enumerate}

\vspace{0.3cm}

% Riquadro Professionale sulla Relazione Equazioni-Dimensioni
\begin{tcolorbox}[colback=blue!5!white, colframe=blue!75!black, title=\textbf{Rettifica}]
Nello spazio vettoriale $\mathbb{R}^n$, ogni equazione cartesiana lineare indipendente agisce come un vincolo che rimuove un grado di libertà.

\begin{itemize}
    \item \textbf{Regola della Dimensione:}
    \[ \dim(U) = n - \text{numero di equazioni linearmente indipendenti} \]
    Dove $n$ è la dimensione dello spazio (es. 4 per $\mathbb{R}^4$).
    
    \item \textbf{Interpretazione Matriciale:}
    Un sistema di $k$ equazioni in $n$ incognite corrisponde a una matrice $k \times n$. Ogni equazione è una \textbf{RIGA} della matrice.
    \[ \dim(U) = n - \text{rango(Matrice dei coefficienti)} \]
\end{itemize}
\end{tcolorbox}

\vspace{0.5cm}

\subsection*{5.2 Intersezione di Sottospazi ($U \cap W$)}
il metodo di risoluzione dipende da come sono forniti i dati.

\textbf{CASO A: Equazioni + Equazioni}
Se $U$ e $W$ sono entrambi definiti da equazioni cartesiane, l'intersezione deve soddisfare \textbf{tutte} le equazioni contemporaneamente.
\begin{enumerate}
    \item Si scrive un unico sistema contenente le equazioni di $U$ e quelle di $W$.
    \item Si risolve il sistema omogeneo risultante.
    \item La soluzione è la base di $U \cap W$.
\end{enumerate}

\textbf{Esempio $\mathbb{R}^4$:}
$U: \begin{cases} x_1+x_2=0 \\ x_3=0 \end{cases}, \quad W: \begin{cases} x_1-x_4=0 \\ x_2+x_3=0 \end{cases}$
\\ $\implies$ Sistema unico $4 \times 4$. Se la soluzione è solo $(0,0,0,0)$, allora $\dim(U \cap W) = 0$.

\medskip

\textbf{CASO B: Span + Equazioni (Metodo di Sostituzione)}
\textit{: Inserire il "vettore generico" dello Span dentro l'equazione cartesiana.}

\textbf{Esempio Svolto:}
Siano $U: x - y + z = 0$ e $W = \text{span}\{ (1, 0, 1), (0, 1, 0) \} \subset \mathbb{R}^3$.
\begin{enumerate}
    \item \textbf{Costruzione Vettore Generico di $W$:}
    \[ w = \alpha(1, 0, 1) + \beta(0, 1, 0) = (\alpha, \beta, \alpha) \]
    
    \item \textbf{Sostituzione nell'equazione di $U$:}
    L'equazione chiede che $(coord_X) - (coord_Y) + (coord_Z) = 0$. Sostituiamo:
    \[ (\alpha) - (\beta) + (\alpha) = 0 \implies 2\alpha - \beta = 0 \]
    
    \item \textbf{Relazione tra parametri:}
    Dall'equazione ricaviamo il vincolo: $\mathbf{\beta = 2\alpha}$.
    
    \item \textbf{Calcolo Base:}
    Torniamo al vettore generico $(\alpha, \beta, \alpha)$ e applichiamo $\beta = 2\alpha$:
    \[ w_{int} = (\alpha, 2\alpha, \alpha) = \alpha(1, 2, 1) \]
    \textbf{Base Intersezione:} $\{ (1, 2, 1) \}$. \textbf{Dim:} 1.
\end{enumerate}

\vspace{0.5cm}

\textbf{CASO C: Span + Span}
\textit{: Convertire il sottospazio più semplice in equazioni cartesiane e ricadere nel Caso B.}



\textbf{Esempio Svolto:}
Siano $U = \text{span}\{ (1, 1, 0), (0, 0, 1) \}$ e $W = \text{span}\{ (1, 0, 2), (0, 1, 0) \}$.

\begin{enumerate}
    \item \textbf{Conversione di $U$ in Cartesiana:}
    Osservando i generatori di $U$: $x$ e $y$ sono uguali nel primo vettore e nulli nel secondo. $z$ è libero.
    Si nota a occhio che $x = y$. Oppure, risolvendo il sistema, l'equazione di $U$ è:
    \[ x - y = 0 \]
    
    \item \textbf{Vettore Generico di $W$:}
    \[ w = \alpha(1, 0, 2) + \beta(0, 1, 0) = (\alpha, \beta, 2\alpha) \]
    
    \item \textbf{Sostituzione (Ricaduta nel Caso B):}
    Inseriamo le coordinate di $w$ nell'equazione $x - y = 0$:
    \[ (\alpha) - (\beta) = 0 \implies \mathbf{\alpha = \beta} \]
    
    \item \textbf{Calcolo Base:}
    Sostituiamo $\beta = \alpha$ nel vettore generico $(\alpha, \beta, 2\alpha)$:
    \[ w_{int} = (\alpha, \alpha, 2\alpha) = \alpha(1, 1, 2) \]
    \textbf{Base Intersezione:} $\{ (1, 1, 2) \}$. \textbf{Dim:} 1.
\end{enumerate}

\vspace{0.3cm}
\hrule
\vspace{0.3cm}

\subsection*{5.3 Somma di Sottospazi e Formula di Grassmann}
La dimensione della somma $U+W$ si calcola mettendo insieme tutti i generatori di $U$ e $W$ in un'unica matrice e calcolandone il rango (riduzione a scala).

\begin{tcolorbox}[colback=yellow!10!white,colframe=orange!85!black,title=\textbf{FORMULA DI GRASSMANN}]
\[ \dim(U + W) = \dim(U) + \dim(W) - \dim(U \cap W) \]

\textbf{Somma Diretta ($\oplus$):}
Si dice che la somma è diretta ($U \oplus W$) se e solo se l'intersezione è banale:
\[ \dim(U \cap W) = 0 \]
In questo caso: $\dim(U \oplus W) = \dim(U) + \dim(W)$.
\end{tcolorbox}
\vspace{0.5cm}

\subsection*{Come determinare la Somma ($U+W$)}
La somma richiede i \textbf{generatori}. L'obiettivo è unire le basi di $U$ e $W$ e scartare i vettori ridondanti.

\begin{tcolorbox}[colback=green!5!white,colframe=green!50!black,title=\textbf{Procedimento}]
Per calcolare $U+W$ bisogna avere \textbf{solo generatori}.
\begin{enumerate}
    \item Se hai delle equazioni, \textbf{risolvile} prima per trovare una base (converti in Span).
    \item Metti \textbf{tutti} i vettori (di $U$ e di $W$) come righe di un'unica matrice.
    \item Riduci a scala con Gauss (Rango).
    \item Le righe non nulle formano una base di $U+W$.
\end{enumerate}
\end{tcolorbox}

\textbf{CASO A: Span + Span (caso ideale)}
abbiamo tutto quel che ci serve.
\begin{itemize}
    \item Costruiramo la matrice unendo i generatori di $U$ e quelli di $W$.
    \item Riduciamo a scala. Il numero di righe non nulle è $\dim(U+W)$.
\end{itemize}

\textbf{CASO B: Span + Equazioni (Il caso misto)}
Convertiamo l'equazione in generatori prima di procedere.
\begin{enumerate}
    \item Risolviamo l'equazione del sottospazio cartesiano (trova le variabili libere) per ottenere la sua base.
    \item Ora ricadiamo nel Caso A: uniamo tutti i vettori in una matrice e riduciamo.
\end{enumerate}

\textbf{CASO C: Equazioni + Equazioni}
Non possiamo sommare equazioni direttamente. Dobbiamo prima convertire i due sottospazi in Span.
\begin{enumerate}
    \item Risolviamo separatamente il sistema di $U$ e il sistema di $W$ per trovare le rispettive basi.
    \item Uniamo le due basi in una matrice e riduciamo (Caso A).
\end{enumerate}

\medskip

\textbf{Esempio Svolto:}
$U = \text{span}\{ (1, 1, 0) \}$ e $W: x - y + z = 0$ in $\mathbb{R}^3$.

\begin{enumerate}
    \item \textbf{Convertire $W$ in Span:}
    $y = x + z$. Variabili libere $x, z$.
    Base $W$: $\{ (1,1,0), (0,1,1) \}$.
    
    \item \textbf{Unione Generatori:}
    Consideriamo insieme $u_1$ e i nuovi $w_1, w_2$:
    $\{ (1,1,0)_U, (1,1,0)_W, (0,1,1)_W \}$.
    
    \item \textbf{Gauss:}
    Il vettore $(1,1,0)$ è ripetuto. La matrice ridotta avrà rango 2.
    \\ $\dim(U+W) = 2$. Base: $\{(1,1,0), (0,1,1)\}$.
\end{enumerate}

\section*{6. Spazi di Matrici e Dimensioni}
\textit{Come calcolare la dimensione di sottospazi definiti da condizioni matriciali (simmetria, prodotti, traccia).}

\begin{tcolorbox}[colback=purple!5!white,colframe=purple!75!black,title=\textbf{Gradi di Libertà}]
La dimensione di uno spazio di matrici è il numero di parametri "liberi" che puoi scegliere.
\[ \dim(X) = \text{Dimensione Totale} - \text{Numero di Vincoli Indipendenti} \]
Dove la \textit{Dimensione Totale} di $M(m,n)$ è $m \times n$.
\end{tcolorbox}

\subsection*{6.1 Vincoli Standard}
Esempi di come i vincoli più comuni riducono la dimensione in $M(n,n)$:

\begin{itemize}
    \item \textbf{Matrici Simmetriche ($A=A^T$):}
    Si scelgono solo la diagonale e il triangolo superiore.
    \[ \dim(\text{Sym}_n) = \frac{n(n+1)}{2} \]
    \textit{Esempio $n=3$:} Dim = 6.
    
    \item \textbf{Matrici Antisimmetriche ($A=-A^T$):}
    \textbf{La diagonale è obbligatoriamente 0 nelle matrici antisimmetriche}. Si sceglie solo il triangolo superiore.
    \[ \dim(\text{Skew}_n) = \frac{n(n-1)}{2} \]
    \textit{Esempio $n=3$:} Dim = 3.
    
    \item \textbf{Traccia Nulla ($\text{tr}(A)=0$):}
    Impone 1 sola equazione sulla diagonale ($a_{11} + \dots + a_{nn} = 0$).
    Toglie esattamente \textbf{1} grado di libertà.
\end{itemize}

\vspace{0.4cm}
\hrule
\vspace{0.4cm}

\subsection*{6.2 Esercizi Svolti (tratti dall'esame)}

\textbf{Esercizio A: Vincoli su righe e colonne}
\\ Sia $X = \{ A \in M(4,4, \mathbb{R}) : A e_1 = 0 \land A^T e_1 = 0 \}$. Calcolare $\dim(X)$.

\textit{1. Vincoli:}
\begin{itemize}
    \item $\mathbf{A e_1 = 0}$:
    Moltiplicare per $e_1$ estrae la prima colonna.
    $\implies$ \textbf{La 1ª colonna deve essere zero.}
    
    \item $\mathbf{A^T e_1 = 0}$:
    Moltiplicare la trasposta per $e_1$ estrae la prima colonna di $A^T$. Ma la prima colonna della trasposta corrisponde esattamente alla prima riga della matrice originale.
    $\implies$ \textbf{La 1ª riga deve essere zero.}
\end{itemize}

\textit{2. Rappresentiamola:}
I vincoli non eliminano interamente le colonne 2, 3 e 4, ma "tagliano via" la prima riga.
Ecco cosa succede alla matrice:

\[
A = \begin{pmatrix}
\mathbf{0}_{11} & \mathbf{0}_{12} & \mathbf{0}_{13} & \mathbf{0}_{14} \\
\mathbf{0}_{21} & \color{blue}* & \color{blue}* & \color{blue}* \\
\mathbf{0}_{31} & \color{blue}* & \color{blue}* & \color{blue}* \\
\mathbf{0}_{41} & \color{blue}* & \color{blue}* & \color{blue}*
\end{pmatrix}
\begin{matrix}
\leftarrow \text{Vincolo } A^T e_1 = 0 \text{ (Riga 1 annullata)} \\
\\
\\
\\
\end{matrix}
\]
\[
\uparrow
\]
\[
\text{Vincolo } A e_1 = 0 \text{ (Colonna 1 annullata)}
\]

\textit{3. Calcolo della Dimensione:}
Le caselle segnate con gli zeri neri sono bloccate dai vincoli.
Le caselle segnate con gli asterischi blu ($*$) rimangono \textbf{variabili libere}.
Si forma un sottomatrice $3 \times 3$ di elementi liberi.

\[ \dim(X) = 3 \times 3 = \mathbf{9} \]

\textbf{Esercizio B: Simmetria + Traccia}
\\ Sia $X = \{ A \in M(3, \mathbb{R}) : A = A^T, \text{tr}(A) = 0 \}$. Calcolare $\dim(X)$.

\textit{Calcolo sequenziale:}
\begin{enumerate}
    \item \textbf{Spazio Totale:} $M(3,3)$ ha dimensione 9.
    \item \textbf{Vincolo Simmetria:} Riduce la dimensione a $\frac{3(4)}{2} = 6$.
    (Parametri liberi: 3 sulla diagonale + 3 sopra la diagonale).
    \item \textbf{Vincolo Traccia:} Aggiunge l'equazione $a_{11} + a_{22} + a_{33} = 0$.
    Questa equazione "toglie" 1 grado di libertà (ad esempio, $a_{33}$ non è più libero, ma dipende dagli altri due).
\end{enumerate}

\[ \dim(X) = 6 (\text{simmetriche}) - 1 (\text{traccia}) = \mathbf{5} \]
\vspace{0.4cm}
\section*{7. Sistemi Lineari e Teorema di Rouché-Capelli}


\hrule
\vspace{0.4cm}

\subsection*{7.1 Definizioni Fondamentali}
Dato un sistema lineare di $m$ equazioni in $n$ incognite:
\[ Ax = b \]
Definiamo due matrici chiave:
\begin{enumerate}
    \item \textbf{Matrice Incompleta ($A$):} La matrice dei soli coefficienti ($m \times n$).
    \item \textbf{Matrice Completa ($A|b$):} La matrice ottenuta aggiungendo la colonna dei termini noti ($m \times (n+1)$).
\end{enumerate}

\vspace{0.3cm}

\begin{tcolorbox}[colback=red!5!white,colframe=red!75!black,title=\textbf{TEOREMA DI ROUCHÉ-CAPELLI}]
Il sistema ha soluzioni se e solo se il rango della matrice incompleta è uguale al rango della matrice completa.

\[ \text{rk}(A) = \text{rk}(A|b) \]

\textbf{Analisi dei Casi:}
\begin{enumerate}
    \item \textbf{CASO IMPOSSIBILE (0 soluzioni):}
    \[ \text{rk}(A) < \text{rk}(A|b) \]
    Il sistema è incompatibile. I termini noti creano una contraddizione (es. $0=1$).
    
    \item \textbf{CASO POSSIBILE (Soluzioni esistenti):}
    \[ \text{rk}(A) = \text{rk}(A|b) = r \]
    Il sistema ammette soluzioni. Per sapere quante, confrontiamo il rango comune $r$ con il numero di incognite $n$.
\end{enumerate}
\end{tcolorbox}

\subsection*{7.2 Conteggio delle Soluzioni}
Una volta stabilito che il sistema è risolubile ($\text{rk}(A) = \text{rk}(A|b) = r$), distinguiamo due sottocasi:

\textbf{A. Soluzione Unica ($r = n$)}
Se il rango è uguale al numero delle incognite, non ci sono variabili libere.
Il sistema ha esattamente \textbf{1 soluzione}.
\textit{(Tipico dei sistemi quadrati con determinante diverso da zero).}

\textbf{B. Infinite Soluzioni ($r < n$)}
Se il rango è minore delle incognite, ci sono variabili che non sono vincolate (variabili libere).
\[ \text{Numero di parametri liberi} = n - r \]
Si dice che il sistema ammette $\infty^{n-r}$ soluzioni.

\vspace{0.3cm}
\begin{tcolorbox}[colback=yellow!10!white,colframe=orange!85!black,title=\textbf{Esempio Pratico}]
Sistema in 3 incognite ($n=3$). Calcoliamo i ranghi.
\begin{itemize}
    \item Se $\text{rk}(A)=2$ e $\text{rk}(A|b)=3 \implies$ \textbf{0 soluzioni}.
    \item Se $\text{rk}(A)=3$ e $\text{rk}(A|b)=3 \implies$ \textbf{1 soluzione} ($n=r$).
    \item Se $\text{rk}(A)=2$ e $\text{rk}(A|b)=2 \implies$ \textbf{$\infty^1$ soluzioni} ($n-r = 3-2 = 1$ parametro).
\end{itemize}
\end{tcolorbox}

\vspace{0.4cm}
\hrule
\vspace{0.4cm}

\subsection*{7.3 Sistemi Omogenei ($Ax = 0$)}
Un sistema è omogeneo se la colonna dei termini noti è nulla ($b=0$).
\begin{itemize}
    \item \textbf{Proprietà:} Un sistema omogeneo è \textbf{sempre compatibile}, perché $\text{rk}(A)$ è sempre uguale a $\text{rk}(A|0)$.
    \item Ammette sempre almeno la soluzione banale $x = (0,0,\dots,0)$.
    \item Se $r < n$, ammette infinite soluzioni (autovettori).
\end{itemize}

\subsection*{7.4 Esercizi Svolti}

\textbf{Esercizio 1:}
\[ \begin{cases} x + y = 1 \\ x + y = 2 \end{cases} \]
Matrice $A = \begin{pmatrix} 1 & 1 \\ 1 & 1 \end{pmatrix}$, $A|b = \begin{pmatrix} 1 & 1 & | & 1 \\ 1 & 1 & | & 2 \end{pmatrix}$.
\begin{itemize}
    \item $\text{rk}(A) = 1$ (righe identiche).
    \item $\text{rk}(A|b) = 2$ (la sottomatrice $\begin{pmatrix} 1 & 1 \\ 1 & 2 \end{pmatrix}$ ha det $\neq 0$).
\end{itemize}
Poiché $1 \neq 2$, il sistema è \textbf{impossibile}.

\vspace{0.3cm}

\textbf{Esercizio 2: Infinite Soluzioni}
\[ x + y - z = 0 \quad (\text{1 equazione, 3 incognite}) \]
\begin{itemize}
    \item $n = 3$.
    \item $\text{rk}(A) = 1$ (c'è un solo pivot).
    \item $b=0$, quindi $\text{rk}(A|b)=1$. Compatibile.
\end{itemize}
Soluzioni: $\infty^{n-r} = \infty^{3-1} = \infty^2$.
Ci sono 2 parametri liberi (es. $y$ e $z$). Soluzione: $x = -y + z$.

\vspace{0.3cm}

\textbf{Esercizio 3: Parametrico}
Discutere al variare di $k$: $\begin{cases} kx + y = 1 \\ x + y = 2 \end{cases}$
\\ Determinante $A = k - 1$.
\begin{itemize}
    \item \textbf{Se $k \neq 1$:} Det $\neq 0 \implies \text{rk}(A)=2$. Anche $\text{rk}(A|b)=2$.
    $r=n=2$. \textbf{1 Soluzione unica}.
    
    \item \textbf{Se $k = 1$:} Il sistema diventa $x+y=1, x+y=2$.
    $\text{rk}(A)=1$, $\text{rk}(A|b)=2$. \textbf{Impossibile}.
\end{itemize}

\section*{8. Autovalori e Diagonalizzazione}
\textit{L'obiettivo è trovare una base in cui la matrice diventa semplicissima (diagonale). Questo è possibile solo se la matrice possiede "abbastanza" autovettori indipendenti.}

\hrule
\vspace{0.4cm}

\subsection*{8.1 Definizioni Chiave}
Sia $A \in M(n, \mathbb{K})$. Un vettore $v \neq 0$ è un \textbf{autovettore} relativo all' \textbf{autovalore} $\lambda$ se:
\[ A v = \lambda v \]
Geometricamente, l'azione di $A$ su $v$ è un semplice allungamento/accorciamento di fattore $\lambda$. La direzione non cambia.

\textbf{Matrice Diagonalizzabile:}
$A$ è diagonalizzabile se esiste una matrice invertibile $P$ e una diagonale $D$ tali che:
\[ D = P^{-1} A P \]
\begin{itemize}
    \item $D$ contiene gli autovalori sulla diagonale.
    \item $P$ contiene gli autovettori sulle colonne (nello stesso ordine degli autovalori).
\end{itemize}

\vspace{0.3cm}

\subsection*{8.2 Algoritmo di Diagonalizzazione}

\textbf{Passo 1: Polinomio Caratteristico ($p_A(\lambda)$)}
Calcolare il determinante della matrice $A - \lambda I$:
\[ p_A(\lambda) = \det(A - \lambda I) \]
Le radici del polinomio sono gli \textbf{Autovalori}.

\textbf{Passo 2: Molteplicità Algebrica ($m_a$)}
Per ogni autovalore $\lambda_0$, contare quante volte annulla il polinomio.
Es: se $p(\lambda) = (\lambda - 2)^3 (\lambda + 1)$, allora:
\begin{itemize}
    \item $\lambda = 2 \implies m_a(2) = 3$
    \item $\lambda = -1 \implies m_a(-1) = 1$
\end{itemize}

\vspace{0.2cm}
\begin{tcolorbox}[colback=green!5!white,colframe=green!50!black,title=\textbf{Autovalori Distinti}]
Se la matrice $n \times n$ ha \textbf{$n$ autovalori tutti distinti}, allora è \textbf{automaticamente diagonalizzabile}.
\textit{Motivo:} Se sono distinti, ogni $m_a = 1$. Poiché $1 \le m_g \le m_a$, allora per forza $m_g = 1$. La condizione è soddisfatta.
\end{tcolorbox}
\vspace{0.2cm}

\textbf{Passo 3: Molteplicità Geometrica ($m_g$)}
Se invece ci sono autovalori ripetuti ($m_a > 1$), bisogna verificare che abbiano abbastanza autovettori.
Per ogni autovalore con $m_a > 1$, calcolare la dimensione dell'autospazio:
\[ m_g(\lambda) = n - \text{rk}(A - \lambda I) \]

\vspace{0.3cm}

\begin{tcolorbox}[colback=blue!5!white,colframe=blue!75!black,title=\textbf{CRITERIO DI DIAGONALIZZABILITÀ}]
Una matrice è diagonalizzabile se e solo se per ogni autovalore $\lambda$:
\[ m_g(\lambda) = m_a(\lambda) \]
Se anche solo per un $\lambda$ si ha $m_g < m_a$, la matrice \textbf{non} è diagonalizzabile.
\end{tcolorbox}

\vspace{0.4cm}
\hrule
\vspace{0.4cm}

\subsection*{8.3 Esempio Svolto (Caso Parametrico)}
Sia $A = \begin{pmatrix} 1 & 0 & 1 \\ 0 & 1 & 0 \\ 1 & 0 & 1 \end{pmatrix}$.

\textbf{1. Polinomio Caratteristico:}
\[ \det(A - \lambda I) = \det \begin{pmatrix} 1-\lambda & 0 & 1 \\ 0 & 1-\lambda & 0 \\ 1 & 0 & 1-\lambda \end{pmatrix} \]
Sviluppando lungo la seconda riga (che ha molti zeri):
\[ = (1-\lambda) \cdot [ (1-\lambda)^2 - 1 ] = (1-\lambda)(1 - 2\lambda + \lambda^2 - 1) \]
\[ = (1-\lambda)(\lambda^2 - 2\lambda) = \lambda(1-\lambda)(\lambda - 2) \]
Radici (Autovalori): $\lambda_1 = 0, \quad \lambda_2 = 1, \quad \lambda_3 = 2$.

\textbf{2. Verifica Molteplicità:}
Poiché abbiamo 3 autovalori distinti per una matrice $3 \times 3$, ognuno ha necessariamente $m_a = 1$.
Dato che $1 \le m_g \le m_a$, allora obbligatoriamente $m_g = 1$ per tutti.
\textbf{Conclusione:} La matrice è diagonalizzabile.

\textbf{3. Costruzione di $D$:}
\[ D = \begin{pmatrix} 0 & 0 & 0 \\ 0 & 1 & 0 \\ 0 & 0 & 2 \end{pmatrix} \]

\medskip

\textbf{Nota Tecnica sugli Autovettori:}
Se avessimo dovuto trovare la base diagonalizzante $P$:
\begin{itemize}
    \item Per $\lambda=0$: Risolvere sistema omogeneo su $A - 0I$.
    \item Per $\lambda=1$: Risolvere sistema omogeneo su $A - 1I$.
    \item Per $\lambda=2$: Risolvere sistema omogeneo su $A - 2I$.
\end{itemize}

\vspace{0.5cm}

\begin{tcolorbox}[colback=orange!5!white,colframe=orange!75!black,title=\textbf{PROPRIETÀ FONDAMENTALI E TRU}]
Sia $A$ una matrice $n \times n$ con autovalori $\lambda_1, \dots, \lambda_n$. Valgono sempre le seguenti pr
àèoprietà:

\begin{enumerate}
    \item \textbf{Traccia e Determinante:}
    La somma degli autovalori è la traccia, il prodotto è il determinante.
    \[ \sum_{i=1}^n \lambda_i = \text{tr}(A) \qquad \prod_{i=1}^n \lambda_i = \det(A) \]
    
    \item \textbf{Invertibilità:}
    Una matrice è invertibile se e solo se $\det(A) \neq 0$, ovvero se \textbf{0 non è un autovalore}.
    
    \item \textbf{Matrici Triangolari:}
    Se $A$ è triangolare (superiore o inferiore) o diagonale, gli autovalori sono esattamente gli elementi sulla diagonale principale.
    
    \item \textbf{Relazione con il Rango:}
    Se $\text{rk}(A) = r < n$, allora la dimensione del nucleo è $n-r$.
    Questo significa che l'autovalore $\lambda = 0$ esiste e ha molteplicità geometrica:
    \[ m_g(0) = n - r \]
    \textit{(In parole povere: se il rango scende di $k$, hai almeno $k$ autovalori nulli).}
    
    \item \textbf{Trasformazioni della Matrice:}
    Se $v$ è autovettore di $A$ relativo a $\lambda$:
    \begin{itemize}
        \item \textbf{Potenza:} $A^k$ ha autovalore $\lambda^k$.
        \item \textbf{Shift:} $A + cI$ ha autovalore $\lambda + c$.
    \end{itemize}
    \textit{Esempio:} Se $A$ ha $\lambda=3$, allora $A+5I$ ha autovalore $3+5=8$.
    
    \item \textbf{Invarianti per Similitudine:}
    Due matrici simili (ovvero che rappresentano la stessa applicazione lineare in basi diverse) condividono:
    \textbf{Determinante, Traccia, Rango, Polinomio Caratteristico e Spettro (Autovalori).}
\end{enumerate}
\end{tcolorbox}
\vspace{0.3cm}

\begin{tcolorbox}[colback=blue!5!white,colframe=blue!60!black,title=\textbf{Rettifica}]
Quando diagonalizziamo una matrice ($D = P^{-1} A P$), non stiamo cambiando la trasformazione, ma solo il sistema di riferimento.
Pertanto, la matrice originale $A$ e la matrice diagonale $D$ sono \textbf{matrici simili}.

\textbf{Cosa condividono $A$ e $D$?}
Tutte le proprietà che non dipendono dalla base scelta:
\begin{itemize}
    \item $\det(A) = \det(D) = \text{prodotto degli autovalori}$
    \item $\text{tr}(A) = \text{tr}(D) = \text{somma degli autovalori}$
    \item $\text{rk}(A) = \text{rk}(D)$
    \item Il Polinomio Caratteristico.
\end{itemize}

\textbf{Attenzione:} La matrice $P$ (che contiene gli autovettori) \textbf{non} condivide queste proprietà con $A$. $P$ serve solo a fare la "traduzione" tra le basi.
\end{tcolorbox}

\subsection*{8.4 Esercizi Teorici Svolti sulla diagonalizzazione}
\textit{Spesso nei temi d'esame viene messo uno "stralcio" di dimostrazione riguardante diagonalizzazione,
prodotti scalari e teorema spettrale (che affronteremo in seguito). riguardo alla parte sulla diagonalizzazione
è utile ricordarsi la formula fondamentale e la proprietà $Av = \lambda v$.}

\vspace{0.3cm}

\textbf{Esercizio A: Polinomi di Matrici}
\\ Sia $v$ un autovettore di $A$ relativo all'autovalore $\lambda = 3$.
\\ Calcolare quanto vale il vettore $w = (A^2 - 4A + 2I)v$.

\textit{Svolgimento:}
Sfruttiamo la linearità e la definizione:
\[ w = A^2 v - 4 A v + 2 I v \]
Sappiamo che:
\begin{itemize}
    \item $A v = 3 v$
    \item $A^2 v = 3^2 v = 9 v$ (Proprietà delle potenze)
    \item $I v = v$
\end{itemize}
Sostituendo:
\[ w = 9v - 4(3v) + 2(1v) = (9 - 12 + 2)v = -1 v \]
\textbf{Risultato:} Il vettore risultante è $-v$. (Nota: $v$ è autovettore anche della matrice risultante con autovalore -1).

\vspace{0.4cm}
\hrule
\vspace{0.4cm}

\textbf{Esercizio B: L'Inversa e gli Autovalori}
\\ Dimostrare che se $\lambda \neq 0$ è autovalore di una matrice invertibile $A$ con autovettore $v$, allora $1/\lambda$ è autovalore di $A^{-1}$ con lo stesso autovettore $v$.

\textit{Dimostrazione:}
Partiamo dalla definizione fondamentale:
\[ A v = \lambda v \]
Poiché $A$ è invertibile, esiste $A^{-1}$. Moltiplichiamo entrambi i membri a sinistra per $A^{-1}$:
\[ A^{-1} (A v) = A^{-1} (\lambda v) \]
A sinistra, l'inversa cancella la matrice ($A^{-1}A = I$):
\[ I v = \lambda (A^{-1} v) \]
\[ v = \lambda (A^{-1} v) \]
Poiché $\lambda \neq 0$, possiamo dividere tutto per $\lambda$:
\[ \frac{1}{\lambda} v = A^{-1} v \]
Rileggendola al contrario: $A^{-1} v = \frac{1}{\lambda} v$.
\textbf{C.V.D.} $v$ è autovettore di $A^{-1}$ con autovalore $\lambda^{-1}$.

\vspace{0.4cm}
\hrule
\vspace{0.4cm}

\textbf{Esercizio C: Matrici Nilpotenti}
\\ Una matrice $A$ si dice nilpotente se esiste un intero $k$ tale che $A^k = 0$ (Matrice nulla). Dimostrare che l'unico autovalore possibile per $A$ è $\lambda = 0$.

\textit{Svolgimento:}
Sia $v$ un autovettore (quindi $v \neq 0$) e $\lambda$ il suo autovalore.
\[ A v = \lambda v \]
Applichiamo la matrice $A$ ripetutamente fino ad arrivare alla potenza $k$:
\[ A^k v = \lambda^k v \]
Ma per ipotesi $A^k$ è la matrice nulla ($0$), quindi il primo membro è il vettore nullo:
\[ \vec{0} = \lambda^k v \]
Poiché il vettore $v$ non è nullo (per definizione di autovettore), l'unico modo affinché il prodotto faccia zero è che lo scalare sia nullo:
\[ \lambda^k = 0 \implies \lambda = 0 \]
\textbf{Conclusione:} Tutti gli autovalori di una matrice nilpotente sono 0.

\vspace{0.4cm}
\hrule
\vspace{0.4cm}

\textbf{Esercizio D: Calcolo rapido di Potenze ($A^k$) (tipologia svolta a lezione, prof. Franchetti)}
\\ Sia $A$ una matrice diagonalizzabile tale che $A = P D P^{-1}$.
\\ Trovare una formula per calcolare $A^{100}$ senza fare cento moltiplicazioni.

\textit{Svolgimento:}
Scriviamo la potenza esplicitando la diagonalizzazione:
\[ A^2 = A \cdot A = (P D P^{-1}) \cdot (P D P^{-1}) \]
Grazie all'associatività, i termini centrali si annullano ($P^{-1} P = I$):
\[ A^2 = P D (P^{-1} P) D P^{-1} = P D (I) D P^{-1} = P D^2 P^{-1} \]
Iterando il ragionamento per 100 volte:
\[ A^{100} = P D^{100} P^{-1} \]
\textbf{Perché è utile?} Elevare a potenza una matrice diagonale $D$ è facilissimo: basta elevare a potenza i singoli numeri sulla diagonale.
\[ \text{Se } D = \begin{pmatrix} 2 & 0 \\ 0 & 3 \end{pmatrix} \implies D^{100} = \begin{pmatrix} 2^{100} & 0 \\ 0 & 3^{100} \end{pmatrix} \]
Il calcolo si riduce a due sole moltiplicazioni matriciali finali (per $P$ e $P^{-1}$).

\vspace{0.4cm}
\hrule
\vspace{0.4cm}

\textbf{Esercizio E: Dimostrazione Invarianza Determinante}
\\ Usando la relazione di similitudine $A = P D P^{-1}$, dimostrare formalmente che $\det(A) = \det(D)$.

\textit{Dimostrazione:}
Applichiamo il determinante a entrambi i membri dell'equazione:
\[ \det(A) = \det(P D P^{-1}) \]
Per il \textbf{Teorema di Binet} ($\det(AB) = \det(A)\det(B)$), possiamo spezzare il prodotto:
\[ \det(A) = \det(P) \cdot \det(D) \cdot \det(P^{-1}) \]
Poiché il determinante è uno scalare (un numero), la moltiplicazione è commutativa. Spostiamo i termini:
\[ \det(A) = \det(D) \cdot [\det(P) \cdot \det(P^{-1})] \]
Sappiamo che $\det(P^{-1}) = \frac{1}{\det(P)}$. Quindi il prodotto tra parentesi fa 1:
\[ \det(A) = \det(D) \cdot 1 \]
\[ \det(A) = \det(D) \]
\textbf{Conclusione:} Il determinante della matrice è uguale al prodotto dei suoi autovalori (che sono gli elementi di $D$).
\end{document}